% Copyright (c) 2025 ZKLib Contributors. All rights reserved.
% Released under Apache 2.0 license as described in the file LICENSE.

\section{FRI}\label{sec:fri}

FRI (Fast Reed--Solomon Interactive Oracle Proof of Proximity) tests whether an oracle
$f : \evaldomain \to \field$ is close to a Reed--Solomon code.  The ArkLib formalization
follows the usual two-phase structure: a folding phase, implemented as a sequential composition
of single-round oracle reductions, and a query phase checking round-to-round consistency.  The
main formal design point is that protocol assembly, completeness, and soundness accounting are
expressed through the generic oracle-reduction composition infrastructure from
Chapter~\ref{chap:oracle_reductions}.

\subsection{Protocol specification}\label{sec:fri_spec}

\begin{definition}[FRI input relation]\label{def:fri_input_relation}
\lean{Fri.Spec.inputRelation}
\uses{def:distance_from_code,def:reed_solomon_code}
    The input relation asserts that the initial oracle statement is $\distance$-close to the
    Reed--Solomon code determined by the initial evaluation domain and degree bound.
\end{definition}

\begin{definition}[FRI output relation]\label{def:fri_output_relation}
\lean{Fri.Spec.outputRelation}
\uses{def:fri_input_relation}
    The output relation is the query-round relation on the final folded statement and witness.
    It is the endpoint of both the folding phase and the full FRI reduction.
\end{definition}

\begin{definition}[Folding-phase protocol]\label{def:fri_folding_protocol}
\lean{Fri.Spec.pSpecFold,Fri.Spec.reductionFold}
\uses{def:fri_input_relation,def:fri_output_relation}
    The folding phase sequentially composes the non-final folding reductions and appends the
    final-fold reduction.  In Lean this is the oracle reduction
    \texttt{Fri.Spec.reductionFold}.
\end{definition}

\begin{definition}[Full FRI reduction]\label{def:fri_reduction}
\lean{Fri.Spec.reduction}
\uses{def:fri_folding_protocol}
    The full FRI reduction appends the query-round oracle reduction to the folding phase.  This
    definitional append is exposed by \texttt{Fri.Spec.reduction\_verifier\_eq\_append} on the
    soundness side.
\end{definition}

\subsection{Perfect completeness}\label{sec:fri_completeness}

Completeness is layered.  Per-round honest-protocol obligations are named as bricks; the
folding phase, query phase, and full reduction are then assembled by the generic append and
sequential-composition theorems.

\begin{definition}[Completeness residuals]\label{def:fri_completeness_residual}
\lean{Fri.Spec.Completeness.foldPhasePerfectCompletenessResidual,Fri.Spec.Completeness.queryRoundPerfectCompletenessResidual,Fri.Spec.Completeness.reductionPerfectCompletenessResidual}
\uses{def:fri_reduction}
    These predicates isolate the per-phase completeness obligations consumed by the composed
    FRI completeness theorem.
\end{definition}

\begin{theorem}[FRI perfect completeness]\label{thm:fri_completeness}
\lean{Fri.Spec.Completeness.foldPhase_perfectCompleteness,Fri.Spec.Completeness.queryRound_perfectCompleteness,Fri.Spec.Completeness.reduction_perfectCompleteness}
\uses{def:fri_reduction,def:fri_completeness_residual}
    Assuming the named per-round and per-phase completeness bricks, the composed FRI reduction is
    perfectly complete.  The proof uses the generic composition keystones rather than duplicating
    FRI-specific append reasoning.
\end{theorem}

\subsection{Soundness accounting}\label{sec:fri_soundness}

The FRI soundness file now contains both the error-budget definitions and the compositional
theorems tying those budgets to the verifier, conditional on the expected per-round and
per-phase soundness obligations.

\begin{definition}[Per-round and total error]\label{def:fri_error}
\lean{Fri.Spec.roundError,Fri.Spec.queryRoundError,Fri.Spec.queryError,Fri.Spec.totalError}
\uses{def:proximity_gap}
    $\mathsf{roundError}_i$ is the BCIKS20-style proximity-gap error for folding round $i$.
    $\mathsf{queryRoundError}_i$ is the query-consistency contribution for one query check,
    $\mathsf{queryError}$ sums those contributions over all query checks, and
    $\mathsf{totalError}$ adds the fold-phase and query-phase budgets.
\end{definition}

\begin{lemma}[Budget projections]\label{lem:fri_error_projections}
\lean{Fri.Spec.roundError_sum_le_totalError,Fri.Spec.roundError_le_totalError,Fri.Spec.queryError_le_totalError}
\uses{def:fri_error}
    The summed fold-round budget, each individual fold-round budget, and the whole query budget
    are each bounded by $\mathsf{totalError}$.
\end{lemma}

\begin{theorem}[Composed FRI soundness]\label{thm:fri_soundness_total_error}
\lean{Fri.Spec.reduction_soundness_of_phases,Fri.Spec.reduction_soundness_totalError,Fri.Spec.reduction_soundness_totalError_of_rounds}
\uses{def:fri_reduction,def:fri_error}
    If the folding phase is sound with the sum of the per-round errors and the query phase is
    sound with $\mathsf{queryError}$, then the appended FRI verifier is sound with error
    $\mathsf{totalError}$.  The round-level theorem further assembles this from individual
    fold-round soundness facts, final-fold soundness, query-round soundness, and the named
    sequential-composition and append residuals.
\end{theorem}

\begin{remark}[Status of the FRI soundness theorem]\label{rem:fri_soundness_status}
    The theorem above is a real composition theorem about the FRI verifier and the
    $\mathsf{totalError}$ budget.  It is not an unconditional low-degree-testing theorem:
    callers still supply the per-round folding soundness, query-round soundness, and generic
    composition residuals.
\end{remark}

\subsection{Batched FRI}\label{sec:batched_fri}

Batched FRI reduces a batch of polynomial oracles to the base FRI protocol through a batching
round and a virtual oracle lens.  Its security development is where ArkLib connects the FRI
query phase to BCIKS20 correlated-agreement statements.

\begin{definition}[Batched FRI reduction]\label{def:batched_fri_reduction}
\lean{BatchedFri.Spec.batchOracleReduction,BatchedFri.Spec.batchedFRIOracleLens,BatchedFri.Spec.liftedFRI,BatchedFri.Spec.batchedFRIreduction}
\uses{def:fri_reduction,def:reed_solomon_code}
    The single batched round forms a random linear combination of the batch and feeds it to a
    lifted base FRI reduction.  The oracle lens routes the lifted FRI oracle queries back to the
    batch oracles.
\end{definition}

\begin{definition}[Batched FRI structural lens package]\label{def:batched_fri_lens_package}
\lean{Fri.batchedFRIOracleLensReduction,Fri.batchedFRIOracleLensReduction_holds}
\uses{def:batched_fri_reduction}
    The structural part of the batching-to-base-FRI reduction is proved definitionally: the
    constructed lens and lifted reduction have the advertised shape.
\end{definition}

\begin{theorem}[Batched FRI query soundness adapters]\label{thm:batched_fri_query_soundness}
\lean{Fri.fri_query_soundness_of_queryRoundProbabilityBoundAndBatchedFRIOracleLens,Fri.fri_query_soundness_of_queryRoundProbabilityBoundAndBatchedFRIOracleLensAndRSAffineLine,Fri.fri_query_soundness_of_queryRoundProbabilityBoundAndBatchedFRIOracleLensAndRSCurve}
\uses{def:batched_fri_lens_package,lem:joint_proximity_bridges}
    The query-round soundness adapters combine the probability bound for the query sampling
    experiment, the batched FRI oracle lens, and a correlated-agreement input.  The affine-line
    and polynomial-curve front doors specialize those adapters to the BCIKS20 Reed--Solomon
    inputs.
\end{theorem}

\begin{theorem}[Small-field batched FRI query soundness]\label{thm:batched_fri_small_field}
\lean{Fri.fri_query_soundness_of_card_le_affineLine,Fri.fri_query_soundness_of_card_le_curve,Fri.fri_soundness_of_card_le_curve}
\uses{thm:batched_fri_query_soundness,thm:bciks_rs_gap}
    In small-field or vacuous regimes, ArkLib discharges the BCIKS20 residual hypotheses using
    the in-tree unconditional correlated-agreement instances, yielding concrete batched FRI
    query-soundness consequences.
\end{theorem}
